Let $(a_k)_{k\ge 1}$ be a sequence of positive integers defined recursively as follows.
\[
 a_1:=1,\qquad a_2:=2.
\]
For each $k\ge 3$, let $a_k$ be the least integer $m$ such that
(i) $m>a_{k-1}$, and
(ii) there exist integers $1\le i<j\le k-1$ with
\[
 m=\sum_{t=i}^j a_t.
\]
(Thus $m$ is the sum of at least two consecutive earlier terms of the sequence.)

The problem asks for the asymptotic behaviour of $a_k$ as $k\to\infty$ (e.g. growth rate, density of the set $\{a_k\}$ in $\mathbb N$, and/or the typical size of the gap $a_k-k$).

Edge conventions: $\mathbb N=\{1,2,3,\dots\}$. “Consecutive” means consecutive indices. Sums must involve at least two terms, i.e. $j-i+1\ge 2$.
